\documentclass[12pt]{article}
\usepackage{amsfonts,amsthm,amsmath,amssymb,mathtools}
\usepackage{mathrsfs}
\usepackage{enumitem}
\usepackage{tikz-cd}
\usepackage{geometry}
\usepackage{mlmodern}
\usepackage{xcolor}
\usepackage{pgfplots}

\geometry{letterpaper, portrait, margin=1in}
\pgfplotsset{compat=1.18}

\setcounter{section}{0}

\renewcommand{\thesection}{\S\arabic{section}}
\renewcommand{\thesubsection}{\arabic{section}}
\newcommand{\dd}{\mathop{}\!\mathrm{d}}

\newtheorem{thm}{Theorem}
\newenvironment{theorem}[1]{
	\renewcommand{\thethm}{#1}
	\begin{thm}
		}{\end{thm}}

\newtheorem{lma}{Lemma}
\newenvironment{lemma}[1]{
	\renewcommand{\thelma}{#1}
	\begin{lma}
		}{\end{lma}}

\theoremstyle{definition}
\newtheorem{ex}{}
\newenvironment{exercise}[1]{
	\renewcommand{\theex}{#1}
	\begin{ex}
		}{\end{ex}}

\title{Written Test 2 Rewrite}
\author{Connor Mooneyhan}
\date{November 14, 2025}

\begin{document}

\maketitle

\begin{exercise}{1(a)}
	Suppose that $A \subset \mathbb{R}$ and $|A| < \infty$.
	Prove that if $A$ is Lebesgue measurable then for every $\varepsilon > 0$ there exists a set $G$ that is the union of finitely many disjoint bounded open intervals such that \begin{equation*}
		|A \setminus G| + |G \setminus A| < \varepsilon.
	\end{equation*}
\end{exercise}
\begin{proof}
	Let $A$ be Lebesgue measurable, and first assume that $A$ is bounded.
	Let $\varepsilon > 0$.
	Then we can choose a closed $F$ such that $F \subset A$ and $|A \setminus F| < \varepsilon/2$, and we can also choose an open $U$ such that $A \subset U$ and $|U \setminus A| < \varepsilon/2$.
	Now express $U$ as the union of open intervals, so $U = \bigcup_{n=1}^\infty I_n$.
	Since $F$ is closed and bounded, we have \begin{equation*}
		F \subset G
	\end{equation*}
	for some reindexed finite subcollection $G = \left\lbrace I_n \right\rbrace_{1}^k$.
	Then $|A \setminus G| \leq |A \setminus F| < \varepsilon/2$ and $|G \setminus A| \leq |U \setminus A| \leq \varepsilon/2$, so \begin{equation*}
		|A \setminus G| + |G \setminus A| < \varepsilon
	\end{equation*}
	as desired.

	Now we consider the general case where just $|A| < \infty$.
	Choose $A_N \coloneq A \cap [-N, N]$ such that $|A \setminus A_N| < \varepsilon/3$, and let $G = \bigcup_{n=1}^k I_n$ be such that $|A_N \setminus G| + |G \setminus A_N| < \varepsilon/3$, using the first case.
	Then \begin{align*}
		A \setminus G = A \cap G^c & = (A_N \cap G^c) \cup (A \setminus A_N \cap G^c)          \\
		                           & = (A_n \setminus G) \cup ((A \setminus A_N) \setminus G).
	\end{align*}
	Thus \begin{align*}
		|A \setminus G| & = |A_N \setminus G| + |(A \setminus A_N) \setminus G| \\
		                & = |A_N \setminus G| + \varepsilon / 3.
	\end{align*}
	Also, $|G \setminus A| \leq |G \setminus A_N| < \varepsilon/3$ by monotonicity, so \begin{align*}
		|A \setminus G| + |G \setminus A| & \leq |A_N \setminus G| + |G \setminus A_N| + \frac{2\varepsilon}{3} \\
		                                  & \leq \varepsilon.
	\end{align*}
	Assume that for any $\varepsilon > 0$, there exists a $G \in \bigcup_{n=1}^k I_n$ such that $|A \setminus G| + |G \setminus A| < \varepsilon$.
	Since $|A \setminus G| < \varepsilon$, there exists an open $\widetilde{G} \supset A \setminus G$ where $|A \setminus G| \geq |\widetilde{G}| + \varepsilon$.
	Thus $A \subset G \cup \widetilde{G} = U$.

	Since $U$ is open, $A \subset U$, and \begin{align*}
		|U \setminus A| & \leq |G \setminus A| + |\widetilde{G}|                \\
		                & \leq |G \setminus A| + |A \setminus G| + \varepsilon,
	\end{align*}
	so $A$ is Lebesgue measurable.
\end{proof}

\begin{exercise}{1(b)}
	State and prove Egorov's Theorem.
\end{exercise}
Statement: Supposee $(X, S, \mu)$ is a measure space with $\mu(X) < \infty$.
Suppose $f_1, f_2, \dots$ is a sequence of $S$-measurable functions from $X$ to $\mathbb{R}$ that converges pointwise on $X$ to a function $f : X \to \mathbb{R}$.
Then for every $\varepsilon > 0$, there exists a set $E \in S$ such that $\mu(X \setminus E) < \varepsilon$ and $f_1, f_2, \dots$ converges uniformly to $f$ on $E$.
\begin{proof}
	Let $\varepsilon > 0$, and fix $n \in \mathbb{Z}^+$.
	For $m \in \mathbb{Z}^+$, let \begin{equation*}
		A_{m, n} = \bigcap_{k=m}^\infty \left\lbrace x \in X : |f_k(x) - f(x)| < \frac{1}{n} \right\rbrace.
	\end{equation*}
	Then the definition of pointwise continuity guarantees that \begin{equation*}
		\bigcup_{m=1}^\infty A_{m,n} = X.
	\end{equation*}
	Moreover, $A_{1, n}, A_{2, n}, \dots$ is an increasing sequence of sets, so \begin{equation*}
		\lim_{m \to \infty} \mu(A_{m,n}) = \mu(X).
	\end{equation*}
	Then there exists an $m_n \in \mathbb{Z}^+$ such that \begin{equation*}
		\mu(X) - \mu(A_{m,n}) < \frac{\varepsilon}{2^n}.
	\end{equation*}
	Letting $n$ vary now, let \begin{equation*}
		E = \bigcap_{n = 1}^\infty A_{m_n, n}.
	\end{equation*}
	Then \begin{align*}
		\mu(X \setminus E) & = \mu \left( X \setminus \bigcap_{n=1}^\infty A_{m_n, n} \right)   \\
		                   & = \mu \left( \bigcup_{n=1}^\infty (X \setminus A_{m_n, n}) \right) \\
		                   & \leq \sum_{n=1}^\infty \mu(X \setminus A_{m_n, n})                 \\
		                   & < \varepsilon.
	\end{align*}

	Now to prove that $f_1, f_2, \dots$ converges uniformly to $f$ on $E$, let $\varepsilon' > 0$ and let $n \in \mathbb{Z}^+$ such that $1/n < \varepsilon'$.
	Then $E \subset A_{m_n, n}$, so \begin{equation*}
		|f_k(x) - f(x)| < \frac{1}{n} < \varepsilon'
	\end{equation*}
	for all $k \geq m_n$ and $x \in E$ due to our choice of $m_n$.
\end{proof}

\begin{exercise}{1(c)}
	Give an example, with justification, of a measure space $(\mathbb{Z}^+, 2^{\mathbb{Z}^+}, \mu)$ such that \begin{equation*}
		\left\lbrace \mu(E) : E \subset \mathbb{Z}^+ \right\rbrace = [0, 1].
	\end{equation*}
	You must show that it is a measure space.
\end{exercise}
\begin{proof}
	Consider $\mu$ defined by \begin{equation*}
		\mu(E) = \sum_{n \in E} 2^{-n},
	\end{equation*}
	and define, $\mu(\varnothing) = 0$.
	The standard way to represent the resulting measures will be as the binary representation of some real number.
	Indeed, $\mu(\mathbb{Z}^+) = .\overline{1} = 1$, and for any $x \in (0, 1)$, $\mu(E) = x$ where $E$ is defined as the set containing exactly those $n$ for which the $n^{\text{th}}$ binary digit of $x$ after the period (henceforth called the $n^{\text{th}}$ ``binary place'') is 1.

	Given disjoint $\left\lbrace E_i \right\rbrace_{\mathbb{Z+}}$, $\mu(\bigcup \lbrace E_i\rbrace)$ has $n^{\text{th}}$ binary place 1 exactly when $n \in E_i$ for some $i$.
	Since $n \in E_i$ for at most one $i$, $\sum \mu(E_i)$ has $n^{\text{th}}$ binary place 0 if $n \notin E_i$ for all $i$ and has $n^{\text{th}}$ binary place equal to the sum of one 1 with all the rest of the terms being 0 if $n \in E_i$ for some $i$.
	Thus \begin{equation*}
		\mu\left(\bigcup \lbrace E_i \rbrace\right) = \sum \mu(E_i)
	\end{equation*}
	as desired.
\end{proof}

\begin{exercise}{2(c)}
	State and prove a theorem that justifies the claim \textit{Integrable functions live mostly on sets of finite measure}.
\end{exercise}
Statement: Suppose $(X, S, \mu)$ is a measure space, $g : X \to [0, \infty]$ is $S$-measurable, and $\int g \dd \mu < \infty$.
Then for every $\varepsilon > 0$, there exists $E \in S$ such that $\mu(E) < \infty$ and \begin{equation*}
	\int_{X \setminus E} g \dd \mu < \varepsilon.
\end{equation*}
\begin{proof}
	Suppose $\varepsilon > 0$.
	Let $P$ be an $S$-partition $A_1, \dots, A_m$ of $X$ where \begin{equation*}
		\int g \dd\mu < \mathcal{L}(g, P) + \varepsilon.
	\end{equation*}
	Let $E$ be the union of the $A_i$ satisfying $\inf_{A_i} f > 0$.
	Then $\mu(E) < \infty$ since otherwise we would have $\mathcal{L}(g, P) = \infty$.
	Then \begin{align*}
		\int_{X \setminus E} g \dd\mu & = \int g \dd\mu - \int \chi_E g \dd \mu                       \\
		                              & < (\mathcal{L}(g, P) + \varepsilon) - \mathcal{L}(\chi_Eg, P) \\
		                              & = \varepsilon,
	\end{align*}
	where the last line follows because $\inf_{A_j} f = 0 $ for all $A_j \not\subset E$.
\end{proof}

\begin{exercise}{3(b)}
	State and prove a theorem that justifies one of Littlewood's Three Principles.
	A brief paraphrase of those principles is \begin{enumerate}[label=\roman*.]
		\item Every Lebesgue measurable set is nearly an open set.
		\item Every Lebesgue integrable function is nearly continuous.
		      (Briefly outline the three major steps of this proof and then provide details for one of those steps.)
		\item Every pointwise convergent sequence of Lebesgue integrable functions is nearly uniformly convergent.
	\end{enumerate}
	If you have already addressed one or more of these principles in previous problems, then make a different choice here.
\end{exercise}
Statement (for ii): Let $(X, S, \mu)$ be a measure space, and let $\varepsilon > 0$.
Then given $f \in \mathcal{L}^1(\mu)$, there exists a continuous function $h$ such that $||f-h||_1 < \varepsilon$.
\begin{proof}
	The proof consists of first proving that we can approximate $\mathcal{L}^1(\mu)$ functions by simple functions, then that we can approximate simple functions by step functions, and finally that we can approximate step functions by continuous functions.
	I'll detail the first step.

	Let $f \in \mathcal{L}^1(\mu)$ and let $\varepsilon > 0$.
	Then there exist simple functions $g_1, g_2 \in \mathcal{L}^1(\mu)$ such that $0 \leq g_1 \leq f^+$ and $0 \leq g_2 \leq f^-$ and \begin{equation*}
		\int(f^+ - g_1)\dd\mu < \frac{\varepsilon}{2} \hspace{1em}\text{and}\hspace{1em} \int(f^- - g_2) \dd\mu < \frac{\varepsilon}{2}.
	\end{equation*}
	Then $g = g_1 - g_2$ is a simple function in $\mathcal{L}^1(\mu)$ and \begin{align*}
		||f - g||_1 & = ||(f^+ - g_1) - (f^- - g_2)||_1                 \\
		            & = \int (f^+ - g_1)\dd\mu + \int (f^- - g_2)\dd\mu \\
		            & < \varepsilon
	\end{align*}
	as desired.
\end{proof}

\begin{exercise}{3(c)}
	Suppose $\lambda$ is Lebesgue measure on $\mathbb{R}$ and $f: \mathbb{R} \to \mathbb{R}$ is a Borel measurable function such that $\int_R |f| \dd \lambda < \infty$.
	Let $g : \mathbb{R} \to \mathbb{R}$ such that \begin{equation*}
		g(x) \coloneq \int_{(x-1, x+1)} f \dd\lambda.
	\end{equation*}
	Prove that $g$ is uniformly continuous on $\mathbb{R}$.
\end{exercise}
\begin{proof}
	I'm very disappointed that I haven't been able to figure this one out in time.
\end{proof}

\end{document}
